%% ma: L12-B14-0044
%% lop: 12
%% chuong: 5
%% bai: 14
%% dang: 12.14.15
%% loai: TN4
%% muc: TH
%% dap_an: "B"
%% nguon: "CĐ2.6.pdf giáo viên gửi 27/10/2026 (M07, Lê Minh Khôi – Nguyễn Anh Khôi, Chương II Chủ đề 6 Ứng dụng vào các bài toán thực tế), Câu 34"
%% kien_thuc: "Bài 14 (tích có hướng), Bài 7 (tọa độ vectơ theo i, j, k)"
%% trang_thai: chinh-thuc
%% ghi_chu: "Giáo viên đã duyệt 27/10/2026"
\begin{ex}
Vectơ lực từ $\vec F$ tác dụng lên một hạt proton có điện tích $q=1{,}6\cdot10^{-19}$ C (Coulomb) chuyển động trong từ trường $\vec B$ với vectơ vận tốc $\vec v$ được cho bởi công thức $\vec F=q\cdot\left[\vec v,\vec B\right]$ (trong đó $\vec v$ tính bằng m/s, $\vec B$ tính bằng Tesla (T), và $\vec F$ tính bằng Newton (N)). Xác định vectơ lực từ tác dụng lên một proton chuyển động trong mặt phẳng $(Oxy)$ với vận tốc $\vec v=10^5\cdot\vec i+10^5\cdot\vec j$ (m/s) đi vào vùng từ trường $\vec B=0{,}3\cdot\vec j$ (T).
\choice{$\vec F=-4{,}8\cdot10^{-15}\cdot\vec k$ (N)}{\True $\vec F=4{,}8\cdot10^{-15}\cdot\vec k$ (N)}{$\vec F=4{,}8\cdot10^{-15}\cdot\vec i$ (N)}{$\vec F=4{,}8\cdot10^{-15}\cdot\vec j$ (N)}
\loigiai{$\left[\vec v,\vec B\right]=\left[\left(10^5;10^5;0\right),(0;0{,}3;0)\right]=\left(0;0;3\cdot10^4\right)$.

$\vec F=1{,}6\cdot10^{-19}\cdot3\cdot10^4\vec k=4{,}8\cdot10^{-15}\vec k$ (N).}
\end{ex}
